% ============================================================
%  PLAN DE MANAGEMENT PROJET (PMP) — Conforme PMBOK 7e éd.
%  Time'Eats — Cellule PMO
%  Projet : Refonte Web
% ============================================================
\documentclass[11pt,a4paper,oneside]{report}
\usepackage{../style/consulting}
\hypersetup{pdftitle={Plan de Management Projet (PMBOK) — Refonte Web — Time'Eats PMO}}
\pagestyle{fancy}\fancyhf{}
\renewcommand{\headrulewidth}{0pt}
\fancyhead[L]{\begin{tikzpicture}[remember picture,overlay]
  \draw[cnNavy,line width=0.5pt](0,0)--(\textwidth,0);
\end{tikzpicture}\vspace{2pt}\timeeatslogo\hspace{4pt}\small\textcolor{cnSlate}{Time'Eats \textbf{·} Plan de Management Projet (PMP) — \textit{Refonte Web}}}
\fancyhead[R]{\small\textcolor{cnSlate}{\thepage}}
\fancyfoot[C]{\tiny\textcolor{cnSlate}{Document confidentiel — Cellule PMO — CESI MAALSI 2026 — Conforme PMBOK 7e éd.}}

\newcommand{\ragV}{\cellcolor{cnGreen!20}\textbf{\textcolor{cnGreen}{Vert}}}
\newcommand{\ragA}{\cellcolor{cnOrange!20}\textbf{\textcolor{cnOrange}{Ambre}}}
\newcommand{\ragR}{\cellcolor{cnRed!20}\textbf{\textcolor{cnRed}{Rouge}}}

\begin{document}\small

\begin{center}
{\LARGE\bfseries\color{cnNavy} Plan de Management Projet}\\[4pt]
\textcolor{cnOrange}{\rule{10cm}{1pt}}\\[4pt]
{\large\color{cnSlate} Phase 2 — Planification \quad|\quad Projet : \textbf{Refonte Web}}\\[2pt]
{\footnotesize\color{cnSlate}\textit{Structure conforme au PMBOK 7\textsuperscript{e} édition — PMI}}
\end{center}

\vspace{4pt}
\renewcommand{\arraystretch}{1.4}
\begin{tabularx}{\textwidth}{L{3.5cm} Y L{3.5cm} Y}
\toprule
\rowcolor{cnNavy}\th{Champ} & \th{Valeur} & \th{Champ} & \th{Valeur} \\
\midrule
Nom du projet   & Refonte Web            & Version PMP    & 1.0 \\
\rowcolor{cnIce}
Chef de projet  & S. Martin              & Date           & 22/04/2026 \\
Sponsor         & DSI Time'Eats          & ESN partenaire & Sopra Steria (forfait 400 j/h) \\
\rowcolor{cnIce}
Budget total    & 300 K€                 & Durée          & 6 mois (Avr–Sep 2026) \\
Approche        & Hybride Agile / jalonnée & Statut       & En cours (avancement 8\%) \\
\bottomrule
\end{tabularx}

\vspace{4pt}
\renewcommand{\arraystretch}{1.3}
\begin{tabularx}{\textwidth}{L{2.5cm} Y}
\toprule
\rowcolor{cnNavy}\th{Historique} & \th{} \\
\midrule
v0.1 — 10/04/2026  & Première trame — cadrage initial par le CP \\
\rowcolor{cnIce}
v0.5 — 18/04/2026  & Revue interne DSI (budget, jalons, risques) \\
v1.0 — 22/04/2026  & Version de référence validée en COPIL d'avril 2026 \\
\bottomrule
\end{tabularx}

% ── 0. Introduction & cycle de vie ────────────────────────
\section*{0 — Introduction, cycle de vie et approche de développement}

Le présent Plan de Management Projet (PMP) décrit la manière dont le projet
\textbf{Refonte Web} est planifié, exécuté, contrôlé et clôturé. Il intègre les plans
subsidiaires et les lignes de base exigés par le \textit{PMBOK Guide} 7\textsuperscript{e}
édition (PMI, 2021), adaptés au contexte DSI de Time'Eats.

\textbf{Cycle de vie :} phases PMBOK \textit{Initiation → Planification → Exécution →
Surveillance \& Maîtrise → Clôture}.

\textbf{Approche de développement :} \textit{hybride}. Un socle \textit{predictive} sécurise
les jalons contractuels avec l'ESN (J0 à Jn) tandis qu'une exécution \textit{adaptive}
(Scrum, sprints de 2 semaines, incrément démontrable) guide la construction produit. Des
revues de performance mensuelles (COPIL) assurent la synchronisation entre les deux rythmes.

\textbf{Destinataires du document :} DG, DSI (sponsor), chef de projet, équipe projet,
Sopra Steria, PMO. \textbf{Cadence de mise à jour :} mensuelle ou à chaque événement majeur
(changement de scope, dépassement d'un seuil d'alerte, révision d'une ligne de base).

% ── 1. Enjeux et objectifs ────────────────────────────────
\section*{1 — Enjeux stratégiques et objectifs SMART}

\subsection*{1.1 Contexte et enjeux}

Time'Eats amorce en 2026 une transformation digitale dont la \textbf{refonte web} est la
pierre angulaire. Le site actuel est daté, peu responsive, non interfaçable avec le futur
CRM, et ne soutient pas la croissance nationale visée. Les enjeux portés par ce projet sont
les suivants :

\renewcommand{\arraystretch}{1.35}
\begin{tabularx}{\textwidth}{L{3.5cm} Y}
\toprule
\rowcolor{cnNavy}\th{Enjeu} & \th{Description} \\
\midrule
Stratégique          & Soutenir la croissance nationale en offrant un canal digital
moderne, socle commun avec le CRM et les applications mobiles clients / fournisseurs. \\
\rowcolor{cnIce}
Commercial           & Générer des leads qualifiés, améliorer le taux de conversion et
différencier Time'Eats sur un marché concurrentiel (restauration collective / traiteurs). \\
Expérience client    & Proposer une expérience fluide, responsive, accessible, cohérente
avec la promesse de marque et le futur espace client. \\
\rowcolor{cnIce}
Opérationnel         & Industrialiser la production web (CMS pour les contributeurs métier,
CI/CD, environnements isolés) et baisser le coût de possession (TCO). \\
Conformité           & Mettre en conformité avec le RGPD, le RGAA 4.1 (accessibilité) et les
standards de sécurité applicatifs (OWASP ASVS niveau 2). \\
\rowcolor{cnIce}
Technologique        & Constituer le socle d'interopérabilité (API-first) pour le CRM, l'App
Mobile Clients et l'App Mobile Fournisseurs. \\
Image \& marque      & Renforcer la notoriété en ligne et aligner l'identité digitale sur la
nouvelle charte graphique Time'Eats. \\
\bottomrule
\end{tabularx}

\subsection*{1.2 Objectifs SMART}

Chaque objectif est formulé selon la méthode \textbf{SMART} (\textit{Spécifique,
Mesurable, Atteignable, Réaliste, Temporellement défini}).

\renewcommand{\arraystretch}{1.3}
\begin{tabularx}{\textwidth}{C{0.7cm} L{2.6cm} Y C{2.1cm} C{2cm}}
\toprule
\rowcolor{cnNavy}\th{\#} & \th{Axe} & \th{Objectif SMART} & \th{Indicateur} & \th{Échéance} \\
\midrule
O1 & Délai              & Livrer la \textbf{v1 en production} conforme au périmètre contractuel                                               & Go-live réalisé       & 30/09/2026 \\
\rowcolor{cnIce}
O2 & Coût               & Tenir le \textbf{budget à $\pm$5\%} des 300 K€ prévus sur l'ensemble du projet                                       & CPI $\geq$ 0,95       & 30/09/2026 \\
O3 & Qualité UX         & Obtenir un score de \textbf{satisfaction MOA $\geq$ 80/100} sur les 4 axes CDQ+F                                      & Enquête post go-live  & 15/10/2026 \\
\rowcolor{cnIce}
O4 & Performance        & Temps de chargement \textbf{$\leq$ 2 s} (p95) sur les pages publiques, score Lighthouse $\geq$ 90                      & Mesure APM / audit    & 20/09/2026 \\
O5 & Accessibilité      & Atteindre la \textbf{conformité RGAA 4.1 niveau AA} ($\geq$ 95\% critères conformes)                                 & Audit accessibilité   & 15/09/2026 \\
\rowcolor{cnIce}
O6 & Sécurité           & Zéro vulnérabilité critique (OWASP ASVS niv.\ 2) au \textbf{pentest pré go-live}                                      & Rapport pentest       & 10/09/2026 \\
O7 & Adoption           & \textbf{+30\% de trafic organique} et taux de rebond $\leq$ 45\% à M+3 du go-live                                     & Google Analytics      & 31/12/2026 \\
\rowcolor{cnIce}
O8 & Intégration        & Intégrer \textbf{100\% des API CRM} spécifiées (contrats v1) au J4                                                   & Tests de contrat verts& 31/08/2026 \\
O9 & Conduite du changt.& \textbf{$\geq$ 95\% des contributeurs} formés au CMS avant la bascule                                                & Taux de formation     & 25/09/2026 \\
\bottomrule
\end{tabularx}

\subsection*{1.3 Facteurs clés de succès (FCS)}

\textit{Disponibilité effective des référents métier (MOA, PO)} ;
\textit{continuité de l'engagement de Sopra Steria et stabilité de l'équipe} ;
\textit{synchronisation rigoureuse avec les projets CRM, Office 365 et Migration Azure} ;
\textit{discipline Agile} (DoR/DoD respectés, backlog priorisé, revues de sprint tenues) ;
\textit{qualité des spécifications API côté CRM} (contrats v1 figés en juin 2026) ;
\textit{sponsoring actif} de la DSI en COPIL et au niveau DG.

% ── 2. Plan de management du contenu (Scope) ──────────────
\section*{2 — Plan de management du contenu (Scope Management Plan)}

\textbf{Énoncé de contenu.} Refonte du site web Time'Eats en application moderne responsive
exposant un \textit{portail public}, un \textit{espace client authentifié}, un \textit{espace
fournisseur}, un \textit{CMS éditorial}, et des intégrations \textit{API CRM} et \textit{SSO
Office 365}, avec conformité \textit{RGAA / RGPD} et pipeline \textit{CI/CD} dédié.

\renewcommand{\arraystretch}{1.3}
\begin{tabularx}{\textwidth}{L{1.5cm} Y}
\toprule
\rowcolor{cnNavy}\th{} & \th{} \\
\midrule
\cellcolor{cnGreen!20}\textbf{In}  & Portail public, espace client, espace fournisseur, CMS,
API CRM (v1), SSO O365, SEO / RGAA / RGPD, pipeline CI/CD, recette fonctionnelle et technique,
transfert en exploitation, formation contributeurs CMS. \\
\rowcolor{cnIce}
\cellcolor{cnRed!20}\textbf{Out} & Applications mobiles clients / fournisseurs (projets
dédiés), déploiement du CRM, Migration Azure (projet dédié), module de paiement en ligne
(phase 2), intégration ERP comptable (hors portefeuille 2026). \\
\bottomrule
\end{tabularx}

\textbf{Processus de contrôle du scope :} le backlog produit est priorisé par le PO, les
éléments entrants sont validés en revue de sprint, et toute modification hors backlog suit
le processus formel de demande de changement (§11). \textbf{WBS} détaillé dans P2-06
(dictionnaire WBS versionné).

\textbf{Livrables majeurs} : \textit{maquettes UX validées} (J1), \textit{socle technique +
env.\ dev} (J2), \textit{MVP portail public + espace client} (J3), \textit{intégration CRM +
espace fournisseur} (J4), \textit{PV de recette} (J5), \textit{go-live + PV de transfert}
(Jn).

% ── 3. Plan de management des exigences (Requirements) ────
\section*{3 — Plan de management des exigences (Requirements Management Plan)}

\renewcommand{\arraystretch}{1.3}
\begin{tabularx}{\textwidth}{L{4cm} Y}
\toprule
\rowcolor{cnNavy}\th{Activité} & \th{Modalité} \\
\midrule
Collecte          & Ateliers métier, interviews MOA (marketing, commerce, fournisseurs), analyse du site existant, benchmark concurrentiel, enquêtes utilisateurs \\
\rowcolor{cnIce}
Formalisation     & User stories (INVEST), critères d'acceptation, DoR / DoD, épopées regroupant les sous-périmètres \\
Priorisation      & \textbf{MoSCoW} (Must / Should / Could / Won't) pondéré par valeur métier × coût ; arbitrage PO hebdomadaire \\
\rowcolor{cnIce}
Traçabilité       & Matrice de traçabilité exigences $\leftrightarrow$ user stories $\leftrightarrow$ tests (outillée Jira + plugin traça.) \\
Validation        & Revue de sprint (PO + MOA), recette fonctionnelle (PV P5-19), recette technique (PV P5-20) \\
\rowcolor{cnIce}
Gestion du changt.& Toute modification d'exigence impactant le baseline passe par P3-14 \\
\bottomrule
\end{tabularx}

% ── 4. Plan de management des délais (Schedule) ───────────
\section*{4 — Plan de management des délais (Schedule Management Plan)}

\textbf{Unité de mesure :} jours ouvrés / sprints de 2 semaines. \textbf{Seuil d'alerte :}
dérive de J+5 ouvrés sur un jalon contractuel $\Rightarrow$ rapport d'écarts (P4-17).
\textbf{Méthode :} chemin critique sur les jalons, \textit{burndown} sprint, suivi
hebdomadaire de la vélocité.

\renewcommand{\arraystretch}{1.3}
\begin{tabularx}{\textwidth}{C{1cm} L{5cm} Y C{2.5cm} C{2cm}}
\toprule
\rowcolor{cnNavy}\th{\#} & \th{Jalon} & \th{Livrable} & \th{Date cible} & \th{Statut} \\
\midrule
J0 & Démarrage officiel                  & Charte signée + kick-off ESN           & 01/04/2026 & \ragV \\
\rowcolor{cnIce}
J1 & Cadrage UX / maquettes validées     & Maquettes UX Sprint 1 validées MOA     & 14/04/2026 & \ragV \\
J2 & Socle technique + env.\ dev         & Env.\ dev accessible, pipeline CI/CD   & 15/05/2026 & \ragA \\
\rowcolor{cnIce}
J3 & MVP portail public + espace client  & Incrément démontrable                  & 15/07/2026 & \ragA \\
J4 & Intégration CRM + esp.\ fournisseur & API CRM branchées, parcours recetté    & 31/08/2026 & \ragA \\
\rowcolor{cnIce}
J5 & Recette fonctionnelle \& technique  & PV recette P5-19 + P5-20 signés        & 20/09/2026 & \ragA \\
Jn & Mise en production \& clôture       & Go-live, PV transfert, REX             & 30/09/2026 & \ragA \\
\bottomrule
\end{tabularx}

\textit{Planning Gantt détaillé : cf.\ P2-07. La ligne de base ordonnancement (schedule
baseline) est figée en COPIL d'avril 2026 et ne peut être modifiée que par demande de
changement approuvée.}

% ── 5. Plan de management des coûts (Cost) ────────────────
\section*{5 — Plan de management des coûts (Cost Management Plan)}

\textbf{Devise :} EUR. \textbf{Méthode d'estimation :} bottom-up sur WBS + forfait ESN +
benchmark historique DSI. \textbf{Contrôle :} Earned Value Management (\textbf{EVM}) mensuel,
indicateurs CPI, SPI, EAC, ETC. \textbf{Seuil d'alerte :} écart > 10\% sur consommé vs
planifié $\Rightarrow$ rapport d'écarts (P4-17) et arbitrage COPIL.

\renewcommand{\arraystretch}{1.3}
\begin{tabularx}{\textwidth}{L{5cm} C{2.2cm} C{2cm} Y}
\toprule
\rowcolor{cnNavy}\th{Poste} & \th{Montant (K€)} & \th{\% total} & \th{Commentaire} \\
\midrule
Intégration Sopra Steria (forfait)   & 165 & 55\% & Dév, intégration CRM, UX/UI — 400 j/h, jalons de paiement J2–J5 \\
\rowcolor{cnIce}
Licences \& outillage (CMS, APM)     &  30 & 10\% & CMS headless, licences dev, outillage QA / APM \\
Ressources internes (PO, TL, 2 devs, MOA) &  45 & 15\% & S. Martin + Tech Lead + 2 devs internes + équipe métier (valorisation JH interne) \\
\rowcolor{cnIce}
Recette \& tests                     &  18 &  6\% & Recette MOA + tests sécurité / charge / pentest \\
Formation \& conduite du changement  &  12 &  4\% & Formation contributeurs CMS + support go-live \\
\rowcolor{cnIce}
Réserve pour aléas (\textit{contingency}) &  30 & 10\% & Sous contrôle du CP — risques identifiés \\
\midrule
\rowcolor{cnNavy!10}\textbf{TOTAL (cost baseline)} & \textbf{300} & \textbf{100\%} & Consommé à date : 18 K€ (6\%) \\
\bottomrule
\end{tabularx}

\textit{Une \textbf{réserve de management} (\textit{management reserve}, non incluse dans
la cost baseline, arbitrée par la DSI) couvre les risques non identifiés. Son utilisation
requiert une demande de changement formelle.}

% ── 6. Plan de management de la qualité (Quality) ─────────
\section*{6 — Plan de management de la qualité (Quality Management Plan)}

\textbf{Politique qualité :} conformité PMBOK + référentiel DSI Time'Eats. \textbf{Standards
applicables :} RGAA 4.1 niveau AA, RGPD, OWASP ASVS niveau 2, charte graphique Time'Eats,
ISO 25010 (qualité logicielle).

\renewcommand{\arraystretch}{1.3}
\begin{tabularx}{\textwidth}{L{4.5cm} Y}
\toprule
\rowcolor{cnNavy}\th{Activité qualité} & \th{Modalité} \\
\midrule
Assurance qualité (QA)   & Revue de code systématique (2 pairs), pipeline CI/CD bloquant sur tests, audit accessibilité en continu, \textit{quality gates} SonarQube \\
\rowcolor{cnIce}
Contrôle qualité (QC)    & Tests unitaires $\geq$ 80\% couverture, tests d'intégration API CRM, tests de charge (objectif p95 $\leq$ 2 s), pentest avant J5 \\
Critères d'acceptation   & DoD partagée : tests verts, doc à jour, RGAA OK, revue UX OK, audit sécurité intermédiaire OK \\
\rowcolor{cnIce}
Métriques                & Vélocité, taux de défauts post-sprint, couverture tests, SLA disponibilité, Lighthouse, satisfaction MOA \\
Non-conformités          & Journal d'anomalies P3-15, traitement prioritaire selon criticité (P1 < 24 h, P2 < 72 h, P3 < 1 sprint) \\
\rowcolor{cnIce}
Amélioration continue    & Rétrospective de sprint, plan d'action par sprint, REX consolidé en clôture (P5-23) \\
\bottomrule
\end{tabularx}

% ── 7. Plan de management des ressources (Resource) ──────
\section*{7 — Plan de management des ressources (Resource Management Plan)}

\textbf{Matrice RACI détaillée : cf.\ P1-04.} L'équipe est structurée en \textit{feature
team} Agile, avec un PO temps plein (CP), un Tech Lead, des développeurs ESN et des
référents métier / sécurité en renfort ponctuel.

\renewcommand{\arraystretch}{1.3}
\begin{tabularx}{\textwidth}{L{4cm} L{3cm} Y C{2cm}}
\toprule
\rowcolor{cnNavy}\th{Ressource} & \th{Rôle} & \th{Activités principales} & \th{\% affec.} \\
\midrule
S. Martin             & Chef de projet / PO    & Pilotage, backlog, COPIL                & 80\% \\
\rowcolor{cnIce}
Référent MOA métier   & Product Owner métier   & Besoin, recette, validation UX          & 40\% \\
Tech Lead interne     & Archi \& intégration   & Archi cible, revue code, intégration    & 60\% \\
\rowcolor{cnIce}
\textbf{Dev interne \#1} & \textbf{Développeur full-stack} & Dév front/back en binôme avec ESN, intégration continue, contribution au \textit{design system} & \textbf{80\%} \\
\textbf{Dev interne \#2} & \textbf{Développeur back / intégration} & Dév API, intégration CRM \& SSO O365, tests automatisés, \textit{pair programming} avec ESN & \textbf{60\%} \\
\rowcolor{cnIce}
QA / Recette interne  & Testeur                & Plans de tests, recette, non-régression & 40\% \\
Sopra Steria          & Intégrateur (forfait)  & Dév front/back, UX/UI, DevOps           & 400 j/h \\
\rowcolor{cnIce}
Référent CI/CD        & Support DevOps         & Pipeline, environnements                & 20\% \\
DPO                   & Conformité RGPD        & DPIA, registre des traitements          & Ponctuel \\
\rowcolor{cnIce}
RSSI                  & Sécurité               & Revue d'architecture sécurité, pentest  & Ponctuel \\
Designer UX / UI (ESN)& Conception            & Maquettes, design system                & Inclus forfait \\
\bottomrule
\end{tabularx}

\textbf{Acquisition / développement / libération :} engagements formalisés avec les
managers métier et la DSI ; montée en compétences via pair programming, \textit{dojos}
techniques et ateliers ; libération planifiée à Jn selon le plan de transfert en
exploitation (P5-22). \textbf{Gestion des conflits :} escalade CP $\rightarrow$ DSI.

\textbf{Intégration des développeurs internes.} Deux développeurs internes Time'Eats
rejoignent l'équipe dès J1 (cadrage UX terminé) en \textit{pair programming} avec
Sopra Steria. Objectifs : (i) \textit{capitalisation} de la connaissance technique en
interne (réversibilité), (ii) \textit{accélération} sur les intégrations sensibles (API CRM,
SSO O365), (iii) \textit{continuité} en phase Run post go-live. Impacts : +60 j/h internes
valorisés (intégrés au poste Ressources internes), révision de la DoD (revue croisée
interne + ESN), onboarding technique documenté.

% ── 8. Plan de management des communications ──────────────
\section*{8 — Plan de management des communications (Communications Management Plan)}

\renewcommand{\arraystretch}{1.3}
\begin{tabularx}{\textwidth}{L{3.8cm} L{3cm} Y C{2cm}}
\toprule
\rowcolor{cnNavy}\th{Communication} & \th{Destinataires} & \th{Canal / Format} & \th{Fréquence} \\
\midrule
Daily stand-up        & Équipe + ESN              & Teams / 15 min                      & Quotidien \\
\rowcolor{cnIce}
Revue de sprint       & PO, MOA, équipe, ESN      & Teams + démo incrément              & Bi-hebdo (fin sprint) \\
Rétrospective         & Équipe projet             & Teams / Miro                        & Bi-hebdo (fin sprint) \\
\rowcolor{cnIce}
Planning de sprint    & PO, équipe, ESN           & Teams + Jira                        & Bi-hebdo (début sprint) \\
CR d'avancement       & DSI, Sponsor              & Email + SharePoint (P3-13)          & Hebdomadaire \\
\rowcolor{cnIce}
Tableau de bord PMO   & DSI, DG                   & Cockpit DSI / PDF (P4-16)           & Mensuel \\
COPIL Refonte Web     & DSI, DG, ESN              & Réunion + CR (P4-18)                & Mensuel \\
\rowcolor{cnIce}
Rapport d'écarts      & DSI, Sponsor              & Email + SharePoint (P4-17)          & Si écart \\
Communication go-live & Utilisateurs finaux       & Intranet, email, affichage          & Avant / après Jn \\
\rowcolor{cnIce}
Demande de changement & Comité changement         & SharePoint (P3-14)                  & Sur demande \\
\bottomrule
\end{tabularx}

\textbf{Archivage :} tous les documents projet sont versionnés sur SharePoint PMO et
rattachés au dossier unique du projet (\textit{project folder}).

% ── 9. Plan de management des risques (Risk) ──────────────
\section*{9 — Plan de management des risques (Risk Management Plan)}

\textbf{Méthode :} identification en atelier trimestriel, analyse qualitative Probabilité $\times$
Impact (échelle 1–5), seuil de criticité $\geq$ 12 traité en COPIL, seuil $\geq$ 16 escaladé
DG. \textbf{Registre complet : P2-09.} \textbf{Stratégies de réponse (PMBOK) :} Éviter,
Transférer, Atténuer, Accepter (menaces) — Exploiter, Partager, Améliorer, Accepter
(opportunités).

\renewcommand{\arraystretch}{1.3}
\begin{tabularx}{\textwidth}{C{0.5cm} L{3.8cm} C{1.2cm} C{1.2cm} C{0.9cm} Y}
\toprule
\rowcolor{cnNavy}\th{\#} & \th{Risque} & \th{P (1–5)} & \th{I (1–5)} & \th{Score} & \th{Réponse} \\
\midrule
R1 & Setup env.\ dev Sopra Steria retardé (VPN S14) & 4 & 4 & 16 & \textbf{Atténuer} — escalade DSI Infra, plan de rattrapage, réserve mobilisée \\
\rowcolor{cnIce}
R2 & Dépendance SSO Office 365 non disponible      & 3 & 4 & 12 & \textbf{Atténuer} — sync hebdo, fallback auth locale temporaire \\
R3 & API CRM instables / specs tardives            & 4 & 5 & 20 & \textbf{Atténuer} — contrats figés juin, mocks, tests de contrat \\
\rowcolor{cnIce}
R4 & Hébergement Azure non prêt pour go-live       & 3 & 5 & 15 & \textbf{Transférer} — plan B on-premise + coord.\ Migration Azure \\
R5 & Scope creep en mode Agile                     & 3 & 3 &  9 & \textbf{Atténuer} — backlog PO, revues, gates budget J3/J4 \\
\rowcolor{cnIce}
R6 & Non-conformité RGAA / RGPD                    & 2 & 4 &  8 & \textbf{Atténuer} — audit continu, DPIA amont \\
R7 & Indisponibilité ressources internes (MOA)     & 2 & 3 &  6 & \textbf{Accepter} — backup MOA identifié \\
\rowcolor{cnIce}
R8 & Attrition d'équipe (CP, Tech Lead)            & 2 & 5 & 10 & \textbf{Atténuer} — doc continue, pair programming, backup CP \\
R9 & Cyberattaque pré go-live                      & 2 & 5 & 10 & \textbf{Transférer} — pentest, WAF, plan de réponse incident \\
\rowcolor{cnIce}
O1 & \textit{Opportunité :} réutilisation du \textit{design system} pour les apps mobiles & 3 & 4 & 12 & \textbf{Exploiter} — capitalisation formelle dès J3 \\
\bottomrule
\end{tabularx}

% ── 10. Plan de management des approvisionnements ─────────
\section*{10 — Plan de management des approvisionnements (Procurement Management Plan)}

\renewcommand{\arraystretch}{1.3}
\begin{tabularx}{\textwidth}{L{4cm} Y}
\toprule
\rowcolor{cnNavy}\th{Objet} & \th{Modalité} \\
\midrule
Type de contrat ESN      & Forfait résultat Sopra Steria (\textbf{FFP} — Firm Fixed Price) — 165 K€ / 400 j/h, jalons de paiement indexés sur J2–J5 \\
\rowcolor{cnIce}
Procédure de sélection   & Déjà attribué (ESN cadre DSI) — revue de conformité cadre achats en début de projet \\
Licences \& logiciels    & Achats via centrale DSI — CMS headless, outillage APM, licences dev et QA \\
\rowcolor{cnIce}
Clauses clés             & SLA support, pénalités de retard sur jalons contractuels, clause de réversibilité, propriété du code source Time'Eats, confidentialité, RGPD \\
Suivi fournisseur        & Point hebdomadaire CP / chef de mission ESN, indicateurs j/h consommés vs planifiés, tableau de bord contractuel mensuel \\
\rowcolor{cnIce}
Critères d'acceptation   & PV recette fonctionnelle + technique signés, pentest OK, documentation remise, transfert de connaissance réalisé \\
Clôture contrat          & PV de fin de prestation, enquête satisfaction ESN, archivage du dossier achats \\
\bottomrule
\end{tabularx}

% ── 11. Plan de management des parties prenantes ──────────
\section*{11 — Plan de management des parties prenantes (Stakeholder Engagement Plan)}

\renewcommand{\arraystretch}{1.3}
\begin{tabularx}{\textwidth}{L{3.5cm} C{2cm} C{2.5cm} Y}
\toprule
\rowcolor{cnNavy}\th{Partie prenante} & \th{Pouvoir / Intérêt} & \th{Niveau cible} & \th{Stratégie d'engagement} \\
\midrule
DG Time'Eats          & Fort / Fort     & Soutien            & COPIL mensuel, tableau de bord DG, escalades majeures \\
\rowcolor{cnIce}
DSI (sponsor)         & Fort / Fort     & Leader             & COPIL, arbitrages, CR hebdo, sponsoring visible \\
Direction Commerciale & Moyen / Fort    & Soutien            & Démos de sprint, ateliers besoins, enquête satisfaction \\
\rowcolor{cnIce}
Direction Marketing   & Moyen / Fort    & Soutien            & Alignement charte, contenus, plan de lancement \\
Utilisateurs finaux   & Faible / Fort   & Soutien            & Tests utilisateurs, formation, com go-live \\
\rowcolor{cnIce}
Contributeurs CMS     & Faible / Fort   & Soutien            & Ateliers CMS, documentation, support go-live \\
Sopra Steria          & Fort / Fort     & Soutien            & Gouvernance contractuelle, point hebdo, COPIL \\
\rowcolor{cnIce}
DPO / RSSI            & Fort / Moyen    & Informé            & Validation DPIA, audit sécurité, pentest \\
Chefs projet liés (CRM, O365, Azure, App Mobile) & Moyen / Fort & Collaboratif & Sync inter-projets (portefeuille), gestion des dépendances \\
\rowcolor{cnIce}
Fournisseurs / clients B2B & Faible / Fort & Informé         & Communication de service, FAQ, support \\
\bottomrule
\end{tabularx}

% ── 12. Plans additionnels ────────────────────────────────
\section*{12 — Plans additionnels : Changement, Configuration, Conduite du changement}

\textbf{Conduite du changement (cf.\ P2-11) :} cartographie des impacts utilisateurs, plan
de communication cible (intranet, email, affichage), ateliers de prise en main, formation
des contributeurs CMS (objectif O9), support renforcé go-live + 30 jours.

\textbf{Gestion de la configuration :} versionnage Git (\textit{trunk-based development}),
tags par sprint, environnements dev / staging / prod isolés et reproductibles (IaC),
\textit{Definition of Done} incluant mise à jour de la documentation technique et
utilisateur, politique de sauvegarde / restauration.

\textbf{Gestion des modifications (Change Control) :} toute évolution hors backlog passe
par une \textit{Demande de Changement} (P3-14), évaluée en termes d'impact scope / coût /
délai / qualité / risques, validée en Comité changement (CP + PO + DSI + Sponsor), tracée
dans le registre des DC et intégrée au baseline par décision formelle.

\textbf{Workflow DC :} émission $\rightarrow$ analyse d'impact $\rightarrow$ comité changement
$\rightarrow$ décision (approuvé / refusé / différé) $\rightarrow$ mise à jour baseline /
backlog $\rightarrow$ communication.

% ── 13. Lignes de base et mesure de performance ──────────
\section*{13 — Lignes de base (Baselines) et mesure de performance}

\renewcommand{\arraystretch}{1.3}
\begin{tabularx}{\textwidth}{L{4cm} Y C{3cm}}
\toprule
\rowcolor{cnNavy}\th{Ligne de base} & \th{Description} & \th{Figée au} \\
\midrule
Scope baseline       & Énoncé de contenu + WBS + dictionnaire WBS            & 22/04/2026 \\
\rowcolor{cnIce}
Schedule baseline    & Planning Gantt jalonné (J0–Jn) — P2-07                & 22/04/2026 \\
Cost baseline        & 300 K€ (hors réserve de management) — P2-08           & 22/04/2026 \\
\rowcolor{cnIce}
Performance baseline & Intégration scope + délais + coûts (EVM mensuel)      & 22/04/2026 \\
\bottomrule
\end{tabularx}

\vspace{4pt}
\textbf{Indicateurs de performance (KPI) :}

\renewcommand{\arraystretch}{1.3}
\begin{tabularx}{\textwidth}{L{4cm} C{2.5cm} Y}
\toprule
\rowcolor{cnNavy}\th{KPI} & \th{Cible} & \th{Source / Fréquence} \\
\midrule
CPI (\textit{Cost Perf.\ Index})  & $\geq$ 0,95  & EVM mensuel (P4-16) \\
\rowcolor{cnIce}
SPI (\textit{Schedule Perf.\ Index}) & $\geq$ 0,95  & EVM mensuel (P4-16) \\
Vélocité (pts / sprint)           & Stable $\pm$10\% & Jira / fin de sprint \\
\rowcolor{cnIce}
Taux de défauts post-sprint       & $\leq$ 5\% & Journal P3-15 / sprint \\
Satisfaction MOA (0–100)          & $\geq$ 80   & Enquête mensuelle \\
\rowcolor{cnIce}
Couverture tests                  & $\geq$ 80\% & CI / sprint \\
Disponibilité env.\ (SLA)         & $\geq$ 99\% & Monitoring / mensuel \\
\bottomrule
\end{tabularx}

% ── 13bis. Alimentation du dashboard PMO ─────────────────
\section*{13bis — Alimentation du dashboard PMO (\textit{Cockpit DSI})}

Le projet Refonte Web alimente le \textbf{Cockpit DSI} (dashboard PMO temps réel consolidé
pour l'ensemble du portefeuille). Cette section définit \textit{quelles données remontent,
depuis quelle source, à quelle fréquence et sous quel format}. Le Cockpit est la source
unique de vérité pour le pilotage de portefeuille (COPIL, DG, DSI).

\subsection*{13bis.1 Données remontées}

\renewcommand{\arraystretch}{1.3}
\begin{tabularx}{\textwidth}{L{4cm} L{3.5cm} C{2.2cm} Y}
\toprule
\rowcolor{cnNavy}\th{Donnée} & \th{Source} & \th{Fréquence} & \th{Utilisation dashboard} \\
\midrule
Avancement (\%)                    & Jira (burndown / tâches faites) & Temps réel (API) & Jauge avancement projet \\
\rowcolor{cnIce}
Tâches total / faites              & Jira                         & Temps réel       & Détail avancement, tendance \\
Consommé (K€)                      & ERP achats + ressources       & Hebdomadaire (ETL) & KPI budget, EVM \\
\rowcolor{cnIce}
CPI / SPI                          & Calcul EVM (PMO)              & Mensuel          & Tableau de bord performance \\
Statut RAG exécution               & CP (déclaration CR hebdo)     & Hebdomadaire     & Heatmap portefeuille \\
\rowcolor{cnIce}
Jalons (date cible, statut)        & Planning Gantt (P2-07)        & À chaque MAJ     & Timeline portefeuille \\
Risques ($\geq$ 12)                & Registre P2-09                & Hebdomadaire     & Top risques agrégés \\
\rowcolor{cnIce}
Incidents / anomalies              & Journal P3-15 (Jira)          & Temps réel (API) & Pastille incidents ouverts \\
Points bloquants                   & CP (CR hebdo)                 & Hebdomadaire     & Liste blocage portefeuille \\
\rowcolor{cnIce}
Satisfaction client (CDQ+F)        & Enquête MOA                   & Mensuelle        & Radar satisfaction \\
Satisfaction équipe                & Rétrospective sprint          & Bi-hebdo         & Courbe tendance équipe \\
\rowcolor{cnIce}
Vélocité / taux défauts            & Jira + SonarQube              & Par sprint       & Qualité livrée \\
Dépendances inter-projets          & PMP + P2-07                   & À chaque MAJ     & Graphe de dépendances \\
\rowcolor{cnIce}
Contacts ESN (jours depuis CR)     & CRM PMO                       & Quotidien        & Vigie relation fournisseur \\
Notes patch / livraisons           & CI/CD (pipeline)              & Temps réel       & Fil des livraisons \\
\bottomrule
\end{tabularx}

\subsection*{13bis.2 Flux et responsabilités}

\renewcommand{\arraystretch}{1.3}
\begin{tabularx}{\textwidth}{L{3cm} L{3cm} L{3cm} Y}
\toprule
\rowcolor{cnNavy}\th{Canal} & \th{Émetteur} & \th{Récepteur} & \th{Mécanisme / Fréquence} \\
\midrule
API Jira           & Dev / PO              & Cockpit (backend) & Webhook temps réel $\rightarrow$ collection \textit{projects.tasks} \\
\rowcolor{cnIce}
CR hebdo (P3-13)   & Chef de projet        & PMO               & Formulaire structuré SharePoint $\rightarrow$ ETL nuit \\
Registre risques   & Chef de projet        & PMO               & Fichier Excel versionné (SharePoint) $\rightarrow$ ingestion hebdo \\
\rowcolor{cnIce}
Budget ERP         & Service achats        & PMO               & Export CSV quotidien $\rightarrow$ rapprochement EVM mensuel \\
Pipeline CI/CD     & Dev / Sopra Steria    & Cockpit           & Webhook GitLab CI $\rightarrow$ collection \textit{patches} \\
\rowcolor{cnIce}
Rétrospective      & Scrum Master / équipe & PMO               & Template Miro $\rightarrow$ synthèse saisie dans Cockpit \\
Enquête satisfaction & MOA                 & PMO               & Forms mensuel $\rightarrow$ alimentation collection \textit{satisfaction\_client} \\
\bottomrule
\end{tabularx}

\subsection*{13bis.3 Règles de gouvernance de la donnée}

\begin{itemize}\setlength{\itemsep}{2pt}
  \item \textbf{Responsable de la donnée} : le CP est \textit{data owner} projet ; le PMO
  est \textit{data steward} portefeuille.
  \item \textbf{Fraîcheur} : une donnée non mise à jour au-delà de son cycle (ex.\ CR hebdo
  > 10 j) génère une alerte automatique au CP et une pastille \textit{ambre} sur le
  Cockpit.
  \item \textbf{Qualité} : contrôles de cohérence automatisés (consommé $\leq$ budget,
  avancement $\in [0,100]$, jalons datés), journal des anomalies de données.
  \item \textbf{Confidentialité} : les données projet sont restreintes à la DSI / PMO ;
  l'export DG est limité aux agrégats de portefeuille.
  \item \textbf{Traçabilité} : historisation mensuelle des valeurs pour analyse de
  tendance et rapprochement ligne de base.
  \item \textbf{Sécurité} : authentification SSO O365 sur le Cockpit, journalisation des
  accès, sauvegarde quotidienne.
\end{itemize}

\subsection*{13bis.4 Cycle de reporting}

\textit{Quotidien} : incidents, pipeline CI/CD, tâches Jira. \quad
\textit{Hebdomadaire} : CR d'avancement, statut RAG, risques, points bloquants. \quad
\textit{Bi-hebdomadaire} : vélocité, satisfaction équipe, qualité de code. \quad
\textit{Mensuel} : EVM (CPI/SPI), satisfaction client, tableau de bord COPIL (P4-16). \quad
\textit{Sur événement} : demande de changement, rapport d'écarts (P4-17), jalon atteint ou
dérivé.

% ── 14. Gouvernance et revues ─────────────────────────────
\section*{14 — Gouvernance et revues}

\renewcommand{\arraystretch}{1.3}
\begin{tabularx}{\textwidth}{L{3.8cm} L{3.5cm} Y}
\toprule
\rowcolor{cnNavy}\th{Instance} & \th{Décisions} & \th{Participants} \\
\midrule
COPIL mensuel     & Validation jalons, arbitrage budget / scope, risques majeurs, DC impactant baseline & DSI, DG, Sponsor, CP, ESN \\
\rowcolor{cnIce}
Revue de sprint   & Validation incrément, priorisation du backlog                   & PO, MOA, équipe, ESN \\
Rétrospective     & Amélioration continue, actions correctives                      & Équipe projet \\
\rowcolor{cnIce}
Comité technique  & Arbitrages techniques, conformité archi cible                   & Tech Lead, Archi DSI, ESN \\
Comité changement & Traitement des Demandes de Changement (P3-14)                   & CP, PO, DSI, Sponsor \\
\rowcolor{cnIce}
Revue qualité / sécurité & Audit qualité, DPIA, pentest, conformité RGAA           & QA, DPO, RSSI, CP \\
\bottomrule
\end{tabularx}

% ── 15. Hypothèses, contraintes et dépendances ───────────
\section*{15 — Hypothèses, contraintes et dépendances}

\renewcommand{\arraystretch}{1.3}
\begin{tabularx}{\textwidth}{L{3cm} Y}
\toprule
\rowcolor{cnNavy}\th{Catégorie} & \th{Éléments} \\
\midrule
Hypothèses      & Disponibilité MOA $\geq$ 40\% ; contrats API CRM figés juin 2026 ; hébergement Azure disponible avant J5 ; équipe Sopra Steria stable \\
\rowcolor{cnIce}
Contraintes     & Budget plafond 300 K€ ; go-live 30/09/2026 (fenêtre commerciale Q4) ; conformité RGAA / RGPD obligatoire ; respect de la charte graphique ; gel des déploiements en décembre \\
Dépendances     & Office 365 (SSO, prérequis J3) ; CRM (API, prérequis J4) ; CI/CD (pipeline, prérequis J2) ; Migration Azure (hébergement cible, prérequis Jn) \\
\rowcolor{cnIce}
Exclusions      & Apps mobiles, module de paiement, ERP comptable, CRM lui-même, Migration Azure \\
\bottomrule
\end{tabularx}

% ── 16. Critères d'acceptation et clôture ────────────────
\section*{16 — Critères d'acceptation et stratégie de clôture}

\textbf{Critères d'acceptation globaux} : PV de recette fonctionnelle (P5-19) et technique
(P5-20) signés ; objectifs SMART O1–O6 atteints ; pentest sans vulnérabilité critique ;
audit RGAA $\geq$ 95\% ; documentation complète remise ; équipe exploitation formée.

\textbf{Stratégie de clôture} : rapport de clôture (P5-21), PV de transfert en exploitation
(P5-22), fiche de retour d'expérience (P5-23), archivage du dossier projet,
libération des ressources, enquête de satisfaction MOA et équipe, capitalisation des
livrables réutilisables (\textit{design system}, composants techniques).

\textbf{Post-projet} : suivi des objectifs d'adoption (O7) à M+3 ; revue des bénéfices à
M+6 en COPIL portefeuille ; plan d'évolution phase 2 (module paiement, nouvelles
fonctionnalités) selon retours utilisateurs.

\vspace{4pt}
\renewcommand{\arraystretch}{1.3}
\begin{tabularx}{\textwidth}{L{3.5cm} Y C{2.5cm}}
\toprule
\rowcolor{cnNavy}\th{Livrable de clôture} & \th{Description} & \th{Référence} \\
\midrule
PV recette fonctionnelle & Validation MOA du périmètre fonctionnel livré & P5-19 \\
\rowcolor{cnIce}
PV recette technique     & Validation DSI de la conformité technique (perf, sécurité, archi) & P5-20 \\
Rapport de clôture       & Bilan scope / coûts / délais / qualité ; écarts vs baselines & P5-21 \\
\rowcolor{cnIce}
PV transfert exploitation & Remise à l'équipe Run (runbook, astreinte, monitoring) & P5-22 \\
Fiche REX                & Retour d'expérience consolidé (succès, difficultés, actions) & P5-23 \\
\rowcolor{cnIce}
Dossier projet archivé   & Archivage SharePoint PMO, conservation 10 ans & — \\
\bottomrule
\end{tabularx}

\vspace{2pt}
\textbf{Validation finale :} le présent PMP est approuvé en COPIL d'avril 2026. Toute
modification substantielle (scope, budget, délais, gouvernance) entraîne une nouvelle
version soumise à validation du Sponsor et du Comité de pilotage, et une communication aux
parties prenantes listées au §11.

\newpage
% ── 17. Glossaire et acronymes ────────────────────────────
\section*{17 — Glossaire et acronymes}

\renewcommand{\arraystretch}{1.3}
\begin{tabularx}{\textwidth}{L{3cm} Y}
\toprule
\rowcolor{cnNavy}\th{Terme / Acronyme} & \th{Définition} \\
\midrule
PMBOK              & \textit{Project Management Body of Knowledge} — corpus de bonnes pratiques en management de projet, publié par le PMI (Project Management Institute). \\
\rowcolor{cnIce}
PMO                & \textit{Project Management Office} — cellule de pilotage du portefeuille projets de la DSI. \\
WBS                & \textit{Work Breakdown Structure} — décomposition hiérarchique du contenu du projet (P2-06). \\
\rowcolor{cnIce}
RACI               & Matrice des responsabilités : \textit{Responsible, Accountable, Consulted, Informed} (P1-04). \\
EVM / CPI / SPI    & \textit{Earned Value Management} et ses indicateurs : \textit{Cost Performance Index} (CPI = EV/AC), \textit{Schedule Performance Index} (SPI = EV/PV). \\
\rowcolor{cnIce}
EAC / ETC          & \textit{Estimate At Completion} / \textit{Estimate To Complete} — prévisions budgétaires en fin de projet. \\
SMART              & \textit{Specific, Measurable, Achievable, Realistic, Time-bound} — méthode de formulation d'objectifs. \\
\rowcolor{cnIce}
DoR / DoD          & \textit{Definition of Ready} / \textit{Definition of Done} — critères d'entrée et de sortie d'une user story. \\
PO / MOA / MOE     & \textit{Product Owner} / Maîtrise d'ouvrage / Maîtrise d'œuvre. \\
\rowcolor{cnIce}
DC                 & Demande de Changement (P3-14). \\
DPIA               & \textit{Data Protection Impact Assessment} — analyse d'impact RGPD. \\
\rowcolor{cnIce}
RGAA               & Référentiel Général d'Amélioration de l'Accessibilité (version 4.1, niveau AA visé). \\
RGPD               & Règlement Général sur la Protection des Données (UE 2016/679). \\
\rowcolor{cnIce}
OWASP ASVS         & \textit{Open Web Application Security Project — Application Security Verification Standard}. \\
SSO                & \textit{Single Sign-On} — authentification unique (ici via Office 365). \\
\rowcolor{cnIce}
FFP                & \textit{Firm Fixed Price} — contrat au forfait résultat. \\
APM                & \textit{Application Performance Monitoring}. \\
\rowcolor{cnIce}
IaC                & \textit{Infrastructure as Code} — gestion des environnements par le code. \\
RAG                & Code couleur de statut : \textit{Red, Amber, Green}. \\
\rowcolor{cnIce}
REX                & Retour d'expérience (P5-23). \\
COPIL              & Comité de pilotage. \\
\bottomrule
\end{tabularx}

\vspace{6pt}
\footnotesize\textit{Documents associés du référentiel PMO : P1-01 Charte, P1-02 Note de
cadrage, P1-03 Fiche identification, P1-04 RACI, P2-06 WBS, P2-07 Gantt, P2-08 Budget,
P2-09 Risques, P2-10 Communication, P2-11 Conduite du changement, P3-12 à P3-15 Exécution,
P4-16 à P4-18 Pilotage, P5-19 à P5-23 Clôture. Structure alignée sur le PMBOK Guide
7\textsuperscript{e} éd.\ (PMI, 2021).}

\end{document}
